\documentclass[11pt]{article}
\usepackage{amsmath,amssymb,mathtools}
\usepackage{fontspec}
\usepackage{unicode-math}
\setmainfont{Latin Modern Roman}
\setmathfont{Latin Modern Math}
\usepackage{xurl}
\usepackage[unicode,colorlinks=true,linkcolor=blue,urlcolor=blue,pdftitle={Problem statement}]{hyperref}
\newcommand{\sref}[2]{\hyperlink{source:#1:#2}{[S#2]}}
\newcommand{\cataloguescope}[1]{\paragraph{Mathematical statement.}#1}
\begin{document}
\section{Research Problems in Function Theory — Problem 2.42}\label{statement}

\subsection{Definitions and mathematical statement}

\paragraph{Shared notation.}$\mathbb N=\{1,2,\ldots\}$, and $\mathbb Z,\mathbb Q,\mathbb R,\mathbb C$ denote the integers, rationals, reals and complex numbers. Any other conventions are specified locally.


An asymptotic path is a continuous curve $\gamma:[0,\infty)\to\mathbb C$ with $|\gamma(t)|\to\infty$. An entire function has asymptotic value $a$ along it if $f(\gamma(t))\to a$. Let $n(r,a,f)$ count all $a$-points in $|z|\le r$ with multiplicity, and let $n_1(r,a,f,\gamma)$ count those lying on the trace of $\gamma$, counted once with their zero multiplicity even if the curve visits them repeatedly.

\paragraph{Statement recorded in the supplied source.}
For $l\ge2$, can one construct an entire function $f$, distinct finite asymptotic values $a_1,\ldots,a_l$, and corresponding asymptotic paths $\gamma_1,\ldots,\gamma_l$, with $n(r,a_k,f)>0$ for all sufficiently large $r$ and every $k$, such that
\[\lim_{r\to\infty}\frac{n_1(r,a_k,f,\gamma_k)}{n(r,a_k,f)}=b_k>0
\quad(k=1,\ldots,l)?\]
Can all the constants $b_k$ be taken to equal $1$? \sref{statement}{2}
Update 2.42 reports Barsegyan's constraint $\sum_k b_k\le1$ for entire functions, so the all-one alternative is excluded when $l\ge2$. \sref{statement}{2}
\subsection{Short English statement}
Can a positive fraction of the points taking each of several asymptotic values lie on its asymptotic path, and can each such fraction be one?
\subsection{Sources}
\begin{description}
\item[\hypertarget{source:statement:1}{S1}] ULAM.AI and the UnsolvedMath Contributors, \emph{UnsolvedMath: A Curated Collection of Open Mathematics Problems}, catalogue entry. CC BY 4.0. \url{https://huggingface.co/datasets/ulamai/UnsolvedMath}.
\item[\hypertarget{source:statement:2}{S2}] W. K. Hayman and E. F. Lingham, \emph{Research Problems in Function Theory} (2018), Chapter 2, Problem 2.42 and Update 2.42. \url{https://arxiv.org/abs/1809.07200}.
\end{description}
\label{end:statement}
\end{document}
